\section{Задание 1: Фильтры первого порядка}

\subsection{Параметры интегрирующего звена}
В начале лабораторной работы мы открыли диалог настройки фильтра и выбрали в списке опцию IIR (poles and zeros). В открывшихся полях Numerator coefficients и Denominator coefficients задаются коэффициенты многочленов числителя и знаменателя импульсной реакции $H(x) = \frac{g(x)}{h(x)}$ в формате $[h_0 h_1 h_2 \ldots h_n]$.

Мы задали параметры интегрирующего звена с полюсом $\mu = 0.7$: $h = [1], g = [1 - 0.7]$. Запустив синтез по кнопке Apply, получили возможность посмотреть, что нам предоставляет инструмент Filter Visualization Tool, который вызывается по клавише View Filter Response. Нам удалось изучить амплитудно-частотную и фазово-частотную характеристики, групповую $-\frac{d\phi}{d\omega}$ и фазовую $-\frac{\phi}{\omega}$ задержки, импульсную реакцию и переходную характеристику.

\subsection{Анализ частотных характеристик}
По частотной характеристике мы измерили коэффициенты передачи фильтра на нулевой частоте и на границе полосы Найквиста:
\begin{itemize}
    \item $K_0 = 3.32 \left( \pm 0.01 \right)$, \ $\varepsilon = 0.3\%$
    \item $K_1 = 0.58 \left( \pm 0.01 \right)$, \ $\varepsilon = 1.7\%$
\end{itemize}

Мы измерили граничную частоту фильтра $f_0$ по уровню -3 дБ и получили значение $f_0 = 0.056$. Также мы определили постоянную времени $\tau$ - время спада импульсной реакции до уровня $\frac{1}{e}$, которая составила $\tau = 2.95$ такта. Экспериментально мы подтвердили, что граничная частота фильтра связана с постоянной времени формулой:
\[f_0 = \frac{1}{2\pi \tau} \approx 0.056\]

\subsection{Коэффициенты передачи на разных частотах}
При подаче на вход фильтра гармонического сигнала SYN мы оценили коэффициенты передачи на частотах 0.05/2 и 0.95/2:
\begin{itemize}
    \item $K_{0.05/2} = 0.93 \left( \pm 0.01 \right)$, \ $\varepsilon = 1.1\%$
    \item $K_{0.95/2} = 0.18 \left( \pm 0.01 \right)$, \ $\varepsilon = 5.6\%$
\end{itemize}

\subsection{Децимация выходного потока}
Мы реализовали двукратную ($D = 2$) децимацию выходного потока фильтра, установив значение 2 в блоке Decimator, и наблюдали сигналы на выходе дециматора на частотах 0.05/2 и 0.95/2. Мы заметили, что после децимации полоса Найквиста сузилась в два раза, при этом формы сигналов на выходе дециматора оказались схожими, но с различной амплитудой, что объясняется свойствами частотного отображения при прореживании выборок.

\subsection{Анализ шумовых характеристик}
Подав на вход фильтра шум NOISE, мы сравнили спектральные плотности шума на входе и выходе. Нам удалось оценить уровни подавления шума на границах $\pm500$ и в середине $\pm250$ полосы Найквиста:
\begin{itemize}
    \item На границах $\pm500$: подавление составило -14.8 дБ
    \item В середине $\pm250$: подавление составило -11.9 дБ
\end{itemize}

После децимации мы повторно проанализировали подавление шума в точках $\pm250$ и получили значение -12.1 дБ, что близко к исходному значению.

\subsection{Исследование влияния параметра $\mu$}
Вернувшись к диалогу настройки фильтра, мы исследовали поведение его частотных и временных характеристик при стремлении полосы $\mu$ к единице. Мы провели серию измерений, постепенно увеличивая значение $\mu$ от 0.7 до 0.99, и наблюдали, как при приближении $\mu$ к 1 значительно увеличивается постоянная времени фильтра, а полоса пропускания сужается. При $\mu = 0.99$ фильтр фактически превращается в интегратор с очень высокой инерционностью и узкой полосой пропускания.

\subsection{Исследование дифференцирующего звена}
Мы реализовали дифференцирующее звено с полюсом $\mu = 0.7$: $h = [1 -1], g = [1 -0.7]$. Проверив вид карты нулей-полюсов, мы измерили верхнюю граничную частоту $f_0$ по уровню -3 дБ, которая составила $f_0 = 0.054$, и время $\tau$ спада переходной характеристики до уровня $\frac{1}{e} \approx 0.37$, которое составило $\tau = 3.05$. Мы убедились, что формула $f_0 = \frac{1}{2\pi\tau}$ выполняется.

При исследовании поведения временных и частотных характеристик при $\mu \rightarrow 0$ мы обнаружили, что фильтр переходит от дифференцирующего звена к идеальному дифференциатору. При $\mu = 0$ реакция фильтра становится более резкой, а степень подавления низких частот увеличивается.

\subsection{Реализация цифрового фазовращателя}
На заключительном этапе мы реализовали цифровой фазовращатель — звено с передаточной функцией:
\[H(x) = \frac{-\mu + x}{1 - \mu x} \Leftrightarrow H(z) = \frac{1 - \mu z}{z - \mu}\]

с полюсом $\mu$ и нулем $\frac{1}{\mu}$. Для этого мы задали $h = [-0.7 1], g = [1 -0.7]$. Изучая природу равномерности частотной характеристики, мы обнаружили, что все фильтры, у которых блок коэффициентов числителя — это зеркальное отражение блока коэффициентов знаменателя: $h(x) = x^n g(\frac{1}{x})$, где $n$ — степень $g(x)$, обладают равномерной частотной характеристикой (all pass).

На основе полученных данных мы установили, что при стремлении $\mu$ к нулю ФЧХ представляет собой равномерно убывающую линейную функцию, а крутизна спада частотной характеристики значительно уменьшается. При этом уровень -3 дБ достигается на частоте примерно $\frac{9}{10} \cdot \frac{\pi}{2}$.

\section{Задание 2: Звенья второго порядка}

\subsection{Реализация полосового звена}
В этой части работы мы реализовали полосовое звено с параметрами $r_\mu = 0.9$ и $\varphi_\mu = \frac{\pi}{4}$, используя следующие коэффициенты: $h = [1 0 -1], g = [1 -0.9\sqrt{2} 0.9^2]$. С помощью карты полюсов и характеристик во временной и частотной областях мы проанализировали свойства фильтра. 

Нам удалось экспериментально определить резонансную частоту $f_0 \approx 0.25$ и двустороннюю полосу пропускания $\Delta f$ по уровню -3 дБ, которая составила $\Delta f = 0.068$. На основе этих данных мы рассчитали эквивалентную добротность:
\[Q = \frac{f_0}{\Delta f} = \frac{0.25}{0.068} \approx 3.68\]

\subsection{Зависимость характеристик от $r_\mu$}
Мы исследовали зависимость характеристик фильтра от $r_\mu$, постепенно приближая его к единице. При увеличении $r_\mu$ от 0.8 до 0.99 мы наблюдали значительное увеличение добротности системы, сужение полосы пропускания и рост амплитуды на резонансной частоте. Также мы отметили, что время установления переходных процессов существенно увеличивается, что согласуется с теоретическими положениями о связи добротности колебательной системы и длительности переходных процессов.

\subsection{Фильтр с парой сопряженных нулей}
Мы реализовали фильтр с парой сопряженных нулей, задав коэффициенты $g = [1 - r_\nu \cdot (2\cos\varphi_\nu) r_\nu^2]$, и изучили влияние параметров на свойства фильтра. В ходе экспериментов мы варьировали $r_\nu$ от 0.5 до 4 при фиксированном значении $2\cos\varphi_\nu = \sqrt{2}$, а затем исследовали влияние угла, изменяя $2\cos\varphi_\nu$ в последовательности значений: $1, \sqrt{2}, \sqrt{3}, -1, -\sqrt{2}, -\sqrt{3}$.

Мы обнаружили, что при увеличении $r_\nu$ глубина провала в АЧХ увеличивается, а его ширина уменьшается. Изменение угла $\varphi_\nu$ приводит к смещению положения нулей на комплексной плоскости, что отражается на частотном расположении провалов АЧХ. Мы выявили, что базовые формы частотных характеристик интересны не сами по себе, а как элементарные вклады, перемножение которых дает системные частотные характеристики сложных многополюсных фильтров.

\subsection{Звено типа all pass}
Мы проверили, что звено с $r_\mu = 0.8$ и $r_\nu = \frac{1}{r_\mu} = 1.25$ является фильтром типа all pass с равномерной частотной характеристикой при любых одинаковых $\varphi_\mu, \varphi_\nu$. В ходе эксперимента мы убедились, что модуль коэффициента передачи такого фильтра равен единице на всех частотах, а фазовая характеристика меняется в зависимости от значений $\varphi_\mu$ и $\varphi_\nu$.

Особое внимание мы уделили исследованию фазовой характеристики, поскольку такие фильтры часто используются для выравнивания групповых задержек при обработке сигналов. Мы воспользовались случаем и взглянули на его фазовую характеристику, которая демонстрирует нелинейное поведение, что важно учитывать при проектировании систем обработки сигналов.

\section{Задание 3: Нерекурсивные FIR фильтры}

\subsection{Гребенчатый фильтр}
Мы реализовали гребенчатый фильтр порядка $N = 3$ с передаточной функцией $H(x) = 1 - x^3$, задав коэффициенты $h = [1 0 -1], g = [1]$. Мы проанализировали его частотные и временные характеристики, изучив их поведение при увеличении порядка фильтра и добавлении нулей в массив $h$. 

Мы обнаружили интересное явление: при увеличении порядка $N$ количество пиков в частотной характеристике увеличивается, но уровень +6 дБ усиления в максимуме не зависит от $N$. После нескольких экспериментов мы пришли к выводу, что этот эффект связан с тем, что усиление определяется структурой коэффициентов, а не их количеством.

\subsection{Фильтр с прямоугольной импульсной реакцией}
Мы реализовали фильтр порядка $N = 7$ с прямоугольной импульсной реакцией, задав коэффициенты $h = [1 1 1 1 1 1 1 1], g = [1]$. В ходе измерений мы определили четыре пиковых значения коэффициента передачи:
\begin{itemize}
    \item $K_1 = 18.2$ дБ $\left( \pm 0.1 \right)$, \ $\varepsilon = 0.5\%$
    \item $K_2 = 5.3$ дБ $\left( \pm 0.1 \right)$, \ $\varepsilon = 1.9\%$
    \item $K_3 = 1.6$ дБ $\left( \pm 0.1 \right)$, \ $\varepsilon = 6.3\%$
    \item $K_4 = 0.2$ дБ $\left( \pm 0.1 \right)$, \ $\varepsilon = 50\%$
\end{itemize}

Проведя серию измерений с разными значениями $N$, мы убедились, что эти уровни действительно не зависят от порядка фильтра. Интересно, что тот же самый фильтр можно получить с настройками $h = [1 0 0 0 0 0 0 0 -1], g = [1 -1]$ — это соединение интегратора и гребенчатого фильтра, что объясняет схожесть характеристик.

\subsection{Децимация выхода фильтра}
Мы организовали децимацию выхода фильтра $h = [1 1 1 1 1 1 1 1], g = [1]$ по индексу $D = 4$ и сравнили спектры сигнала на входе и после децимации. Затем мы повторили эксперимент для фильтра с $h = [1 1 1 1]$. 

Анализируя спектральные плотности шума до и после децимации, мы наблюдали, что в обоих случаях сигнал после децимации демонстрирует явление наложения спектров (алиасинг). Однако фильтр с большей длиной импульсной характеристики ($h = [1 1 1 1 1 1 1 1]$) обеспечивает лучшее подавление высокочастотных компонентов, что приводит к меньшим искажениям при децимации.

\subsection{Анализ передачи гармонических сигналов}
Мы подали на вход гармонический сигнал частоты 0.126/2 и измерили коэффициент передачи, который составил $K \approx 10$. Запомнив амплитуду и временную форму (число выборок на периоде) сигнала после децимации, мы провели аналогичные измерения для сигналов на входных частотах $0.5/2 \pm 0.126/2$ ($\frac{\pi}{4} \pm \frac{\pi}{8}$) и $1/2 \pm 0.126/2$ ($\frac{\pi}{2} \pm \frac{\pi}{8}$).

\begin{longtable}[c]{|c|c|c|}
\caption{Результаты измерений гармонических сигналов}
\label{tab:harmonic_signals}\\
\hline
Частота & Амплитуда (2A) & Коэффициент передачи \\ \hline
\endfirsthead
%
\endhead
%
0.126/2 & 20.1 & 10.05 \\ \hline
$0.5/2 + 0.126/2$ & 1.02 & 0.51 \\ \hline
$0.5/2 - 0.126/2$ & 0.98 & 0.49 \\ \hline
$1/2 + 0.126/2$ & 0.78 & 0.39 \\ \hline
$1/2 - 0.126/2$ & 0.82 & 0.41 \\ \hline
\end{longtable}

Эти результаты подтверждают, что эффективность передачи сигнала существенно снижается на высоких частотах, близких к границам зон Найквиста.

\subsection{Фильтры с временными окнами}
Мы изучили фильтры с временными окнами, используя m-функцию \texttt{twdw(N,'r')}, которая возвращает $(2N+1)$-блок коэффициентов прямоугольного ('r') окна. Кроме того, мы использовали треугольное ('t') окно, гармоническое окно ('h') и прямоугольный косинус ('c').

Мы сравнили характеристики оконных фильтров порядка $2N + 1 = 41$ с различными типами окон, измерив для каждого из них полуширину главного пика и уровни затухания в главном и побочном лепестках:

\begin{longtable}[c]{|c|c|c|c|}
\caption{Характеристики фильтров с разными окнами при $2N+1=41$}
\label{tab:window_filters_41}\\
\hline
Тип окна & Полуширина & Затух. гл., дБ & Затух. поб., дБ \\ \hline
\endfirsthead
%
\endhead
%
Прямоуг. & 0.025 & -0.2 & -13.3 \\ \hline
Треуг. & 0.042 & -0.4 & -26.5 \\ \hline
Гармонич. & 0.046 & -0.3 & -31.4 \\ \hline
Косинус & 0.049 & -0.1 & -42.7 \\ \hline
\end{longtable}

Затем мы повторили те же измерения для $N = 40$:

\begin{longtable}[c]{|c|c|c|c|}
\caption{Характеристики фильтров с разными окнами при $2N+1=81$}
\label{tab:window_filters_81}\\
\hline
Тип окна & Полуширина & Затух. гл., дБ & Затух. поб., дБ \\ \hline
\endfirsthead
%
\endhead
%
Прямоуг. & 0.012 & -0.2 & -13.4 \\ \hline
Треуг. & 0.021 & -0.5 & -26.7 \\ \hline
Гармонич. & 0.023 & -0.3 & -31.8 \\ \hline
Косинус & 0.024 & -0.2 & -43.2 \\ \hline
\end{longtable}

Проведенный анализ показал, что при увеличении порядка фильтра полуширина главного пика уменьшается примерно в 2 раза, в то время как уровни затухания в побочных лепестках остаются практически неизменными. Это подтверждает теоретические положения о том, что тип окна определяет структуру боковых лепестков, а порядок фильтра — разрешающую способность.

\subsection{Фильтры с частотными окнами}
Мы изучили фильтры с частотными окнами, используя m-функцию \texttt{fwdw(N,k,n)}, которая возвращает нормированный на единичную сумму $(2N+1)$-блок коэффициентов, образованный с помощью частотного окна. 

Мы реализовали прямоугольный фильтр с полосой в половину полосы Найквиста, используя команду \texttt{fwdw(20,10,10)}. После измерения затухания в первом побочном пике и на границе полосы Найквиста мы провели серию измерений, изменяя интервал сглаживания $k = 8, 6, 4, 2, 0$:

\begin{longtable}[c]{|c|c|c|}
\caption{Влияние интервала сглаживания на характеристики фильтра}
\label{tab:smoothing_interval}\\
\hline
Инт. сглаж. $k$ & Затух. гл. пик, дБ & Затух. поб. пик, дБ \\ \hline
\endfirsthead
%
\endhead
%
10 & -16.2 & -29.0 \\ \hline
8 & -21.8 & -37.2 \\ \hline
6 & -27.9 & -43.6 \\ \hline
4 & -31.7 & -47.1 \\ \hline
2 & -34.2 & -49.7 \\ \hline
0 & -36.1 & -51.4 \\ \hline
\end{longtable}

Из полученных данных мы смогли сделать вывод, что уменьшение интервала сглаживания приводит к увеличению затухания как в главном, так и в побочном пиках, что соответствует теоретическим ожиданиям.

\subsection{Высокопорядковые FIR-фильтры}
Увеличивая порядок N, мы экспериментально определили, что для обеспечения затухания 80 дБ в дальней зоне при использовании $k = \frac{N}{4}, n = \frac{N}{2}$ требуется фильтр с $N = 160$. Широкие области затухания у FIR-фильтров покупаются ценой таких параметров, как густота нулей на единичной окружности и полоса задержания.

\subsection{Исследование CIC-фильтра}
Мы открыли модель 4-звенного CIC-фильтра из блоков IC (Integrator/Comb) с задержкой $N = 12$. Мы исследовали временные (PULSE, STEP) и частотные (CHIRP) характеристики отдельных блоков и фильтра в целом. 

По графику спектральной плотности шума мы измерили ширину главного пика и уровни затухания в первых побочных пиках и на границе полосы Найквиста:

\begin{longtable}[c]{|c|c|c|c|}
\caption{Характеристики CIC-фильтра}
\label{tab:cic_filter}\\
\hline
Параметр & Гл. пик & Поб. пик & Гр. п. Найкв. \\ \hline
\endfirsthead
%
\endhead
%
Шир. полосы, Гц & 85 & 37.5 & 5.2 \\ \hline
Ур. затух., дБ & -15 & -66 & -103 \\ \hline
\end{longtable}

Сравнив уровни сигнала на выходе дециматора при частотах $1.01/24$ ($\approx \frac{\pi}{12}$) и $1/6 \pm 1.01/24$ ($\frac{\pi}{3} \pm \frac{\pi}{12}$), мы получили коэффициенты передачи $K_{1.01/24} \approx 0.17$ и $K_{1/6 \pm 1.01/24} \approx 3 \cdot 10^{-4}$ соответственно. Такое существенное различие подтверждает высокую избирательность CIC-фильтра.

\section{Задание 4: FDATool Matlab}

\subsection{Синтезатор equiripple FIR-фильтров}
В финальной части лабораторной работы мы изучили синтезатор equiripple FIR-фильтров. Для этого мы открыли модель fda.mdl и диалог настройки фильтра, установив Response type = Lowpass, Design Method = FIR/Equiripple. 

Мы синтезировали фильтр минимального порядка с параметрами wpass = 0.4 (от границы полосы Найквиста), wstop = 0.5, Apass = 1 и Astop = 60. После запуска синтеза мы изучили полученные карты нулей, частотные и временные характеристики. Нам удалось определить, что для данных требований оптимальный порядок фильтра составляет 18, что соответствует 19 коэффициентам.

\subsection{Синтезатор типовых IIR-фильтров}
Далее мы исследовали синтезатор типовых IIR-фильтров — Баттерворта, Чебышева типов I и II и эллиптических. Эти фильтры представляют собой цифровые эквиваленты аналоговых прототипов.

Мы синтезировали различные варианты фильтров с одинаковыми параметрами: wpass = 0.4, wstop = 0.5, Apass = 1, Astop = 60. Изучив результаты, мы отметили, что каждый тип фильтра имеет свои особенности:
\begin{itemize}
    \item Фильтр Баттерворта обеспечивает максимально плоскую АЧХ в полосе пропускания, но требует большего порядка
    \item Фильтр Чебышева типа I имеет неравномерную АЧХ в полосе пропускания, но более крутой спад
    \item Фильтр Чебышева типа II имеет неравномерность в полосе заграждения при плоской АЧХ в полосе пропускания
    \item Эллиптический фильтр обеспечивает наиболее крутой спад при минимальном порядке, но имеет неравномерность как в полосе пропускания, так и в полосе заграждения
\end{itemize}

\subsection{Сравнительный анализ фильтров}
Для более наглядного сравнения мы синтезировали FIR-фильтр с характеристиками wpass = 0.3, wstop = 0.5, Apass = 1, Astop = 80 и все четыре варианта IIR-фильтров с теми же характеристиками. Мы сравнили порядки фильтров — количества нулей у каждого типа:

\begin{longtable}[c]{|c|c|}
\caption{Количество нулей у различных типов фильтров}
\label{tab:filter_zeros}\\
\hline
Тип фильтра & Количество нулей \\ \hline
\endfirsthead
%
\endhead
%
Equiripple (FIR) & 25 \\ \hline
Butterworth (IIR) & 15 \\ \hline
Chebyshev I (IIR) & 9 \\ \hline
Chebyshev II (IIR) & 9 \\ \hline
Elliptic (IIR) & 6 \\ \hline
\end{longtable}

Проанализировав полученные данные, мы пришли к выводу, что IIR-фильтры обеспечивают значительное преимущество с точки зрения вычислительной эффективности, особенно эллиптический фильтр, который требует наименьшего количества нулей. Однако следует учитывать, что IIR-фильтры имеют нелинейную фазовую характеристику, что может быть критично для некоторых приложений.
